\documentclass[14pt,a4paper]{extarticle}

% ============================================================
% ПРЕАМБУЛА
% ============================================================

\usepackage[T2A]{fontenc}
\usepackage[utf8]{inputenc}
\usepackage[russian]{babel}

\usepackage{amsmath, amssymb, amsfonts}
\usepackage{geometry}
\usepackage{graphicx}
\usepackage{float}
\usepackage{booktabs}
\usepackage{array}
\usepackage{longtable}
\usepackage{multirow}
\usepackage{caption}
\usepackage{indentfirst}
\usepackage{setspace}
\usepackage{titlesec}
\usepackage{hyperref}
\usepackage{xcolor}

\geometry{
    left=30mm,
    right=15mm,
    top=20mm,
    bottom=20mm
}

\onehalfspacing
\sloppy

\hypersetup{
    colorlinks=true,
    linkcolor=black,
    urlcolor=black,
    citecolor=black
}

\setlength{\parindent}{1.25cm}
\setlength{\parskip}{0pt}

\titleformat{\section}
{\normalfont\bfseries\Large}
{\thesection.}{0.5em}{}

\titleformat{\subsection}
{\normalfont\bfseries\large}
{\thesubsection.}{0.5em}{}

\titleformat{\subsubsection}
{\normalfont\bfseries\normalsize}
{\thesubsubsection.}{0.5em}{}

\captionsetup{
    labelsep=period,
    justification=centering
}

\newenvironment{explanation}
{
    \vspace{0.2cm}
    \noindent\textbf{Пояснение.}
}
{
    \vspace{0.2cm}
}

% ============================================================
% ДОКУМЕНТ
% ============================================================

\begin{document}

% ============================================================
% ТИТУЛЬНЫЙ ЛИСТ
% ============================================================

\begin{titlepage}
    \begin{center}
        \textbf{МИНИСТЕРСТВО НАУКИ И ВЫСШЕГО ОБРАЗОВАНИЯ\\
        РОССИЙСКОЙ ФЕДЕРАЦИИ}\\[0.5cm]

        \textbf{Федеральное государственное автономное образовательное учреждение\\
        высшего образования}\\[0.3cm]

        \textbf{«Национальный исследовательский университет ИТМО»}\\[1.5cm]

        \textbf{Факультет технологий искусственного интеллекта}\\[2cm]

        \textbf{\Large Лабораторная работа}\\[0.5cm]

        по дисциплине\\
        \textbf{«Функциональный анализ»}\\[1cm]

        \textbf{\Large Регрессия по функциональным данным:}\\[0.2cm]
        \textbf{\Large линейные функционалы и метрические методы}\\[0.8cm]

        \textbf{Кейсы №4 и №6}\\[2cm]
    \end{center}

    \begin{flushright}
        Выполнил:\\
        студент группы J3210\\
        \textbf{Тищенко Павел Валерьевич}
    \end{flushright}

    \vfill

    \begin{center}
        Санкт-Петербург\\
        2026
    \end{center}
\end{titlepage}

% ============================================================
% ОГЛАВЛЕНИЕ
% ============================================================

\tableofcontents
\newpage

% ============================================================
% АННОТАЦИЯ
% ============================================================

\section*{Аннотация}
\addcontentsline{toc}{section}{Аннотация}

В работе рассматривается задача регрессии по функциональным данным, когда объект наблюдения задаётся не конечномерным вектором признаков, а функцией одной переменной на фиксированном интервале.

Сначала исходные функции преобразуются в конечномерные векторы признаков с помощью линейных функционалов интегрального вида. На полученных признаках строится линейная регрессионная модель, исследуется влияние выбора системы функционалов, числа признаков и регуляризации на качество прогноза.

Далее те же функциональные признаки рассматриваются как элементы метрического пространства. На их основе реализуются метрические методы регрессии: ядерная регрессия Надарая--Ватсона с фиксированным и переменным окном, а также робастная схема LOWESS.

Цель работы состоит в сравнении параметрического подхода, основанного на линейных функционалах, и непараметрического метрического подхода с точки зрения качества прогноза, устойчивости к шуму и выбросам, а также интерпретируемости модели.

\newpage

% ============================================================
% ВВЕДЕНИЕ
% ============================================================

\section*{Введение}
\addcontentsline{toc}{section}{Введение}

Во многих прикладных задачах объект наблюдения естественно задаётся не конечномерным вектором признаков, а функцией одной переменной. К таким данным относятся временные ряды, сигналы датчиков, профили нагрузки, суточные температурные кривые, ЭКГ-сигналы и другие объекты, наблюдаемые на фиксированном интервале.

В данной работе рассматривается задача регрессии по функциональным данным. Пусть каждому объекту соответствует функция \(x(t)\), заданная на интервале \([a,b]\), а целевой переменной является число \(y \in \mathbb{R}\). Требуется построить модель, позволяющую по функции \(x(t)\) предсказывать значение \(y\).

Первый подход основан на использовании линейных функционалов. Исходная функция преобразуется в конечномерный вектор признаков с помощью интегралов вида
\[
L_{\varphi}(x)=\int_a^b x(t)\varphi(t)\,dt,
\]
где \(\varphi(t)\) — заранее выбранная весовая функция. После такого преобразования задача сводится к обычной регрессии по числовым признакам.

Второй подход основан на метрических методах регрессии. Полученные функциональные признаки рассматриваются как точки в конечномерном метрическом пространстве. Для построения прогноза используются расстояния между объектами, ядерное сглаживание Надарая--Ватсона, переменная ширина окна и робастная схема LOWESS.

Таким образом, кейс 4 и кейс 6 объединяются в единую исследовательскую схему. В кейсе 4 строится конечномерное представление функциональных данных, а в кейсе 6 на этом представлении применяются метрические методы регрессии. Это позволяет сравнить параметрический и непараметрический подходы к одной и той же задаче.

Цель работы — реализовать и сравнить методы регрессии по функциональным данным на основе линейных функционалов и метрического сглаживания.

Для достижения цели необходимо решить следующие задачи:
\begin{enumerate}
    \item задать пространство функциональных данных и способ их дискретного представления;
    \item построить признаки с помощью нескольких систем линейных функционалов;
    \item реализовать линейную регрессию по функциональным признакам;
    \item исследовать влияние выбора функционалов и их количества на качество прогноза;
    \item восстановить и проанализировать весовую функцию линейной модели;
    \item реализовать ridge-регрессию и исследовать влияние регуляризации;
    \item реализовать ядерную регрессию Надарая--Ватсона с фиксированным и переменным окном;
    \item реализовать робастную схему LOWESS;
    \item сравнить качество моделей по метрикам MAE, MSE, RMSE и \(R^2\);
    \item исследовать устойчивость моделей к шуму, выбросам и изменению параметров.
\end{enumerate}

\newpage

% ============================================================
% 1. ПОСТАНОВКА ЗАДАЧИ
% ============================================================

\section{Постановка задачи регрессии по функциональным данным}

\subsection{Функциональные данные}

В классических задачах машинного обучения объект обычно представляется конечномерным вектором признаков:
\[
x_i = (x_{i1}, x_{i2}, \dots, x_{in}) \in \mathbb{R}^n.
\]

Однако во многих задачах объект имеет более сложную природу и задаётся функцией:
\[
x_i(t), \quad t \in [a,b].
\]

Например, если объектом является суточная температура, то \(t\) может обозначать время в течение суток, а \(x_i(t)\) — значение температуры в момент времени \(t\). Если объектом является сигнал датчика, то \(x_i(t)\) описывает изменение сигнала во времени.

В данной работе рассматриваются функции, определённые на отрезке
\[
[a,b]=[0,1].
\]

Требуется по функции \(x_i(t)\) предсказать числовой ответ
\[
y_i \in \mathbb{R}.
\]

Иными словами, необходимо построить отображение
\[
x_i(t) \mapsto y_i.
\]

\begin{explanation}
Главная сложность состоит в том, что функция — это бесконечномерный объект. Даже если она хранится на сетке, математически она описывает целую кривую, а не один обычный вектор. Поэтому сначала нужно получить из функции конечное число признаков.
\end{explanation}

\subsection{Дискретизация функций}

На практике функция \(x_i(t)\) хранится не как непрерывный объект, а как набор значений на сетке:
\[
t_0,t_1,\dots,t_{N-1}.
\]

В работе используется равномерная сетка:
\[
t_j = a + j\Delta t,\quad j=0,1,\dots,N-1,
\]
где
\[
\Delta t = \frac{b-a}{N-1}.
\]

Тогда функция \(x_i(t)\) представляется в виде набора значений:
\[
x_i(t_0), x_i(t_1), \dots, x_i(t_{N-1}).
\]

Такое представление позволяет численно вычислять интегралы и строить функциональные признаки.

\subsection{Связь кейса 4 и кейса 6}

Кейс 4 посвящён построению регрессии по функциональным данным через линейные функционалы. В нём исходная функция \(x_i(t)\) преобразуется в вектор признаков:
\[
z_i = (z_{i1}, z_{i2}, \dots, z_{im}),
\]
где
\[
z_{ik}=\int_0^1 x_i(t)\varphi_k(t)\,dt.
\]

Кейс 6 посвящён метрическим методам регрессии. В нём предполагается, что объекты уже являются векторами:
\[
x_i \in \mathbb{R}^n.
\]

Поэтому в данной работе кейс 6 применяется к признаковому представлению, полученному в кейсе 4. То есть вместо исходных функций в метрические методы подаются векторы \(z_i\).

Общая схема имеет вид:
\[
x_i(t)
\longrightarrow
z_i=
\left(
\int_0^1 x_i(t)\varphi_1(t)\,dt,
\dots,
\int_0^1 x_i(t)\varphi_m(t)\,dt
\right)
\longrightarrow
\hat y_i.
\]

\begin{explanation}
Именно поэтому два кейса можно объединить в один отчёт. Кейс 4 отвечает за переход от функций к признакам, а кейс 6 использует эти признаки как метрическое пространство, где можно считать расстояния между объектами.
\end{explanation}

% ============================================================
% 2. ОПИСАНИЕ ДАННЫХ
% ============================================================

\section{Описание данных}

В качестве основного набора данных используются синтетические функциональные данные. Каждая функция имеет вид:
\[
x_i(t)=a_i\sin(2\pi t)+b_i\cos(2\pi t)+c_it+\varepsilon_i(t),
\quad t\in[0,1],
\]
где \(a_i\), \(b_i\), \(c_i\) — случайные коэффициенты, а \(\varepsilon_i(t)\) — шум функции.

Целевой ответ задаётся формулой:
\[
y_i=
2\int_0^1x_i(t)\sin(2\pi t)\,dt
-
\int_0^1x_i(t)t\,dt
+\eta_i,
\]
где \(\eta_i\) — шум в ответе.

Такой способ генерации данных удобен тем, что истинная зависимость ответа от функции задаётся через линейные функционалы. Поэтому можно проверить, насколько хорошо разные системы функционалов и разные методы регрессии восстанавливают эту зависимость.

\begin{explanation}
Функция \(x_i(t)\) содержит синусоидальную, косинусоидальную и линейную компоненты. Поэтому ожидается, что тригонометрический базис должен хорошо описывать такие данные, так как он напрямую связан с синусами и косинусами.
\end{explanation}

\begin{figure}[H]
    \centering
    \includegraphics[width=0.85\textwidth]{report_outputs/figures/sample_functions.png}
    \caption{Примеры сгенерированных функциональных данных}
    \label{fig:sample_functions}
\end{figure}

% ============================================================
% 3. ЛИНЕЙНЫЕ ФУНКЦИОНАЛЫ
% ============================================================

\section{Построение признаков с помощью линейных функционалов}

\subsection{Понятие линейного функционала}

Функционалом называется отображение, которое каждой функции сопоставляет число. В данной работе используются интегральные функционалы вида
\[
L_{\varphi}(x)=\int_0^1 x(t)\varphi(t)\,dt,
\]
где \(\varphi(t)\) — фиксированная весовая функция.

Если рассматривать пространство \(L_2[0,1]\), то выражение
\[
\int_0^1 x(t)\varphi(t)\,dt
\]
можно понимать как скалярное произведение функций \(x(t)\) и \(\varphi(t)\):
\[
L_{\varphi}(x)=\langle x,\varphi\rangle.
\]

\begin{explanation}
Простыми словами, функционал измеряет, насколько функция \(x(t)\) согласована с некоторым шаблоном \(\varphi(t)\). Если \(\varphi(t)\) — синус, то функционал показывает, насколько сильно в функции \(x(t)\) выражена синусоидальная компонента.
\end{explanation}

\subsection{Проверка линейности функционала}

Функционал \(L_{\varphi}\) называется линейным, если для любых функций \(x(t)\), \(y(t)\) и любых чисел \(\alpha,\beta \in \mathbb{R}\) выполняется:
\[
L_{\varphi}(\alpha x+\beta y)
=
\alpha L_{\varphi}(x)+\beta L_{\varphi}(y).
\]

Проверим это свойство:
\[
L_{\varphi}(\alpha x+\beta y)
=
\int_0^1(\alpha x(t)+\beta y(t))\varphi(t)\,dt.
\]

По линейности интеграла:
\[
L_{\varphi}(\alpha x+\beta y)
=
\alpha\int_0^1x(t)\varphi(t)\,dt
+
\beta\int_0^1y(t)\varphi(t)\,dt.
\]

Следовательно,
\[
L_{\varphi}(\alpha x+\beta y)
=
\alpha L_{\varphi}(x)+\beta L_{\varphi}(y).
\]

Значит, функционал \(L_{\varphi}\) является линейным.

\begin{explanation}
Это свойство важно, потому что линейные функционалы сохраняют линейную структуру пространства функций. Если функция является линейной комбинацией двух других функций, то значение функционала также раскладывается в такую же линейную комбинацию.
\end{explanation}

\subsection{Системы функций для построения признаков}

Для построения признаков выбирается система функций:
\[
\varphi_1(t),\varphi_2(t),\dots,\varphi_m(t).
\]

Тогда для каждой функции \(x_i(t)\) строится вектор:
\[
z_i=(z_{i1},z_{i2},\dots,z_{im}),
\]
где
\[
z_{ik}=\int_0^1 x_i(t)\varphi_k(t)\,dt.
\]

В работе рассматриваются три системы функционалов.

\subsubsection{Индикаторы подотрезков}

Первая система основана на индикаторах подотрезков:
\[
\varphi_k(t)=\mathbf{1}_{[\tau_{k-1},\tau_k]}(t).
\]

Такие функционалы измеряют вклад функции на отдельных участках интервала.

\begin{explanation}
Если интервал \([0,1]\) разбить на несколько частей, то каждый функционал показывает, насколько велика функция на конкретном участке. Это похоже на локальные средние значения функции.
\end{explanation}

\subsubsection{Тригонометрический базис}

Вторая система основана на тригонометрических функциях:
\[
1,\ \sin(2\pi t),\ \cos(2\pi t),\ \sin(4\pi t),\ \cos(4\pi t),\dots
\]

Такие функционалы измеряют периодические компоненты функции.

\begin{explanation}
Тригонометрический базис особенно полезен, если функция содержит волны или периодические колебания. В нашем случае данные специально содержат \(\sin(2\pi t)\) и \(\cos(2\pi t)\), поэтому такой базис должен быть хорошо согласован с природой данных.
\end{explanation}

\subsubsection{Полиномиальный базис}

Третья система основана на степенных функциях:
\[
1,t,t^2,\dots,t^{m-1}.
\]

Такие функционалы измеряют связь функции с гладкими полиномиальными трендами.

\begin{explanation}
Полиномиальный базис хорошо описывает плавные изменения функции. Однако если данные имеют ярко выраженные периодические компоненты, одних полиномов может быть недостаточно.
\end{explanation}

\subsection{Численное вычисление интегралов}

Поскольку функции заданы на дискретной сетке, интегралы вычисляются численно. В работе используется формула трапеций:
\[
\int_0^1 x(t)\varphi_k(t)\,dt
\approx
\sum_{j=0}^{N-2}
\frac{
x(t_j)\varphi_k(t_j)+x(t_{j+1})\varphi_k(t_{j+1})
}{2}
(t_{j+1}-t_j).
\]

Таким образом, для каждой функции \(x_i(t)\) и каждой весовой функции \(\varphi_k(t)\) вычисляется числовой признак \(z_{ik}\).

\begin{explanation}
Метод трапеций заменяет площадь под графиком суммой площадей маленьких трапеций. Чем плотнее сетка, тем точнее приближается интеграл.
\end{explanation}

% ============================================================
% 4. ЛИНЕЙНАЯ РЕГРЕССИЯ
% ============================================================

\section{Линейная регрессия по функциональным признакам}

\subsection{Модель линейной регрессии}

После построения признаков задача регрессии сводится к обычной линейной модели:
\[
\hat y_i=\beta_0+\sum_{k=1}^{m}\beta_k z_{ik}.
\]

Здесь \(z_{ik}\) — значение \(k\)-го линейного функционала на функции \(x_i(t)\), \(\beta_k\) — коэффициент модели, \(\beta_0\) — свободный член.

\begin{explanation}
На этом этапе функция уже заменена набором чисел. Поэтому можно применять обычную линейную регрессию, как в задачах с табличными данными.
\end{explanation}

\subsection{Матричная форма модели}

Обозначим через \(Z\) матрицу функциональных признаков:
\[
Z=
\begin{pmatrix}
z_{11} & z_{12} & \dots & z_{1m}\\
z_{21} & z_{22} & \dots & z_{2m}\\
\vdots & \vdots & \ddots & \vdots\\
z_{n1} & z_{n2} & \dots & z_{nm}
\end{pmatrix}.
\]

Добавим к ней столбец единиц для свободного члена и получим матрицу:
\[
X=
\begin{pmatrix}
1 & z_{11} & z_{12} & \dots & z_{1m}\\
1 & z_{21} & z_{22} & \dots & z_{2m}\\
\vdots & \vdots & \vdots & \ddots & \vdots\\
1 & z_{n1} & z_{n2} & \dots & z_{nm}
\end{pmatrix}.
\]

Тогда модель записывается в виде:
\[
\hat y=X\beta.
\]

\subsection{Метод наименьших квадратов}

Для оценки коэффициентов используется метод наименьших квадратов. Минимизируется функционал:
\[
Q(\beta)=(y-X\beta)^T(y-X\beta).
\]

Этот функционал равен сумме квадратов ошибок:
\[
Q(\beta)=\sum_{i=1}^{n}(y_i-\hat y_i)^2.
\]

Найдём условие минимума. Производная по \(\beta\) равна:
\[
\frac{\partial Q}{\partial \beta}
=
-2X^T(y-X\beta).
\]

Приравнивая производную к нулю, получаем нормальные уравнения:
\[
X^TX\hat\beta=X^Ty.
\]

Если матрица \(X^TX\) обратима, то решение имеет вид:
\[
\hat\beta=(X^TX)^{-1}X^Ty.
\]

\begin{explanation}
Метод наименьших квадратов подбирает такие коэффициенты, чтобы предсказания модели были как можно ближе к истинным значениям. Квадраты ошибок используются для того, чтобы большие ошибки штрафовались сильнее.
\end{explanation}

\subsection{Ridge-регрессия}

При большом числе признаков или при сильной зависимости признаков друг от друга оценки МНК могут стать неустойчивыми. Для повышения устойчивости используется ridge-регрессия:
\[
Q_{\lambda}(\beta)
=
(y-X\beta)^T(y-X\beta)+\lambda\|\beta\|_2^2,
\]
где \(\lambda \geq 0\) — параметр регуляризации.

Нормальные уравнения принимают вид:
\[
(X^TX+\lambda I)\hat\beta_{\lambda}=X^Ty.
\]

Отсюда:
\[
\hat\beta_{\lambda}
=
(X^TX+\lambda I)^{-1}X^Ty.
\]

\begin{explanation}
Регуляризация добавляет штраф за слишком большие коэффициенты. Это помогает избежать ситуации, когда модель слишком сильно подстраивается под обучающие данные и плохо работает на новых объектах.
\end{explanation}

\subsection{Восстановление весовой функции}

Линейную модель можно записать в функциональной форме:
\[
\hat y(x)=\beta_0+\sum_{k=1}^{m}\beta_k\langle x,\varphi_k\rangle.
\]

Введём функцию весов:
\[
w(t)=\sum_{k=1}^{m}\beta_k\varphi_k(t).
\]

Тогда:
\[
\langle x,w\rangle
=
\left\langle x,\sum_{k=1}^{m}\beta_k\varphi_k \right\rangle.
\]

По линейности скалярного произведения:
\[
\left\langle x,\sum_{k=1}^{m}\beta_k\varphi_k \right\rangle
=
\sum_{k=1}^{m}\beta_k\langle x,\varphi_k\rangle.
\]

Следовательно,
\[
\hat y(x)=\beta_0+\langle x,w\rangle.
\]

\begin{explanation}
Функция \(w(t)\) показывает, какие участки исходной функции \(x(t)\) сильнее влияют на прогноз. Если \(w(t)\) велико по модулю на некотором участке, то значения функции \(x(t)\) на этом участке особенно важны для модели.
\end{explanation}

\begin{figure}[H]
    \centering
    \includegraphics[width=0.85\textwidth]{report_outputs/figures/weight_functions_comparison.png}
    \caption{Сравнение восстановленных весовых функций}
    \label{fig:weight_functions_comparison}
\end{figure}

% ============================================================
% 5. ЭКСПЕРИМЕНТЫ ПО ЛИНЕЙНОЙ МОДЕЛИ
% ============================================================

\section{Экспериментальное исследование линейной модели}

\subsection{Сравнение систем функционалов}

Для оценки качества разных систем функционалов были построены признаки на основе индикаторов подотрезков, тригонометрического базиса и полиномиального базиса. Для каждой системы обучалась линейная модель, после чего качество оценивалось на тестовой выборке.

Основными метриками качества являются:
\[
MAE=\frac{1}{n}\sum_{i=1}^{n}|y_i-\hat y_i|,
\]
\[
MSE=\frac{1}{n}\sum_{i=1}^{n}(y_i-\hat y_i)^2,
\]
\[
RMSE=\sqrt{MSE},
\]
\[
R^2=1-\frac{\sum_{i=1}^{n}(y_i-\hat y_i)^2}
{\sum_{i=1}^{n}(y_i-\bar y)^2}.
\]

\begin{explanation}
MAE показывает среднюю абсолютную ошибку. RMSE также измеряет ошибку, но сильнее штрафует большие отклонения. Метрика \(R^2\) показывает, какую долю изменчивости целевой переменной объясняет модель.
\end{explanation}

\begin{figure}[H]
    \centering
    \includegraphics[width=0.85\textwidth]{report_outputs/figures/basis_comparison_rmse.png}
    \caption{Сравнение качества разных систем функционалов}
    \label{fig:basis_comparison_rmse}
\end{figure}

На рисунке \ref{fig:basis_comparison_rmse} сравниваются системы функционалов по ошибке RMSE. Чем меньше значение RMSE, тем точнее модель восстанавливает зависимость между функцией \(x(t)\) и ответом \(y\).

\subsection{Влияние числа функционалов}

Следующий эксперимент связан с изменением числа функционалов \(m\). При малом \(m\) модель может не получать достаточно информации о функции. При слишком большом \(m\) возрастает риск переобучения и неустойчивости коэффициентов.

\begin{figure}[H]
    \centering
    \includegraphics[width=0.85\textwidth]{report_outputs/figures/rmse_vs_m.png}
    \caption{Зависимость ошибки от числа функционалов}
    \label{fig:rmse_vs_m}
\end{figure}

\begin{figure}[H]
    \centering
    \includegraphics[width=0.85\textwidth]{report_outputs/figures/train_test_rmse_vs_m.png}
    \caption{Влияние числа функционалов на обучение и обобщение}
    \label{fig:train_test_rmse_vs_m}
\end{figure}

По графикам можно оценить баланс между недообучением и переобучением. Если ошибка на обучении уменьшается, а ошибка на тесте начинает расти, это указывает на переусложнение модели.

\subsection{Влияние регуляризации}

Для исследования влияния регуляризации обучались ridge-модели при различных значениях параметра \(\lambda\). Малые значения \(\lambda\) почти не отличаются от обычного МНК, а большие значения сильнее уменьшают коэффициенты модели.

\begin{figure}[H]
    \centering
    \includegraphics[width=0.85\textwidth]{report_outputs/figures/ridge_lambda_curve.png}
    \caption{Зависимость качества ridge-регрессии от параметра регуляризации}
    \label{fig:ridge_lambda_curve}
\end{figure}

Регуляризация полезна тогда, когда увеличение числа признаков делает модель неустойчивой. В этом случае штраф на норму коэффициентов позволяет получить более гладкое и устойчивое решение.

\subsection{Анализ прогнозов и ошибок}

Для лучшей линейной модели были построены графики истинных и предсказанных значений, а также распределение остатков.

\begin{figure}[H]
    \centering
    \includegraphics[width=0.8\textwidth]{report_outputs/figures/linear_true_vs_pred.png}
    \caption{Истинные и предсказанные значения для линейной модели}
    \label{fig:linear_true_vs_pred}
\end{figure}

Если точки на рисунке \ref{fig:linear_true_vs_pred} расположены близко к диагонали, значит модель хорошо предсказывает целевую переменную.

\begin{figure}[H]
    \centering
    \includegraphics[width=0.8\textwidth]{report_outputs/figures/linear_residuals_distribution.png}
    \caption{Распределение ошибок линейной модели}
    \label{fig:linear_residuals_distribution}
\end{figure}

Распределение остатков позволяет проверить, нет ли систематического смещения модели. Если ошибки сосредоточены около нуля, модель в среднем не завышает и не занижает прогноз.

% ============================================================
% 6. МЕТРИЧЕСКИЕ МЕТОДЫ
% ============================================================

\section{Метрические методы регрессии}

\subsection{Гипотеза непрерывности}

Метрические методы регрессии основаны на гипотезе непрерывности: близким объектам соответствуют близкие ответы.

В контексте данной работы объектом для метрической регрессии является не исходная функция \(x_i(t)\), а её признаковое представление:
\[
z_i=(z_{i1},z_{i2},\dots,z_{im}).
\]

Таким образом, расстояния считаются между векторами функциональных признаков.

\begin{explanation}
После кейса 4 каждая функция уже превращена в точку в пространстве признаков. Если две функции похожи с точки зрения выбранных функционалов, то расстояние между их векторами \(z_i\) будет малым.
\end{explanation}

\subsection{Метрика в пространстве функциональных признаков}

В качестве расстояния между объектами используется евклидова метрика:
\[
\rho(z,z_i)=
\sqrt{
\sum_{k=1}^{m}(z_k-z_{ik})^2
}.
\]

Здесь \(z\) — вектор признаков нового объекта, а \(z_i\) — вектор признаков объекта обучающей выборки.

\begin{explanation}
Евклидово расстояние измеряет обычную геометрическую близость между точками. Для метрических методов важно масштабировать признаки, так как признаки большого масштаба могут начать доминировать в расстоянии.
\end{explanation}

\subsection{Ядерная регрессия Надарая--Ватсона}

Оценка Надарая--Ватсона строит прогноз как взвешенное среднее ответов обучающих объектов:
\[
a(z;X^\ell,h)=
\frac{
\sum_{i=1}^{\ell} y_i K\left(\frac{\rho(z,z_i)}{h}\right)
}{
\sum_{i=1}^{\ell} K\left(\frac{\rho(z,z_i)}{h}\right)
}.
\]

Здесь \(K(r)\) — ядро, \(h>0\) — ширина окна, \(\rho(z,z_i)\) — расстояние между объектами.

Чем ближе объект \(z_i\) к объекту \(z\), тем больший вес получает его ответ \(y_i\).

\begin{explanation}
Прогноз получается не через одну глобальную формулу, как в линейной регрессии, а через локальное усреднение соседних объектов. Близкие объекты влияют на прогноз сильнее, дальние — слабее.
\end{explanation}

\subsection{Связь с локально взвешенным методом наименьших квадратов}

Оценку Надарая--Ватсона можно получить как решение задачи локально взвешенного метода наименьших квадратов при постоянной локальной модели:
\[
f(z_i)\approx \theta.
\]

Тогда минимизируется функционал:
\[
Q(\theta;z)
=
\sum_{i=1}^{\ell}
K\left(\frac{\rho(z,z_i)}{h}\right)
(\theta-y_i)^2.
\]

Обозначим
\[
w_i=
K\left(\frac{\rho(z,z_i)}{h}\right).
\]

Тогда:
\[
Q(\theta;z)=\sum_{i=1}^{\ell}w_i(\theta-y_i)^2.
\]

Найдём минимум:
\[
\frac{dQ}{d\theta}
=
2\sum_{i=1}^{\ell}w_i(\theta-y_i)=0.
\]

Отсюда:
\[
\theta \sum_{i=1}^{\ell}w_i
=
\sum_{i=1}^{\ell}w_iy_i.
\]

Следовательно,
\[
\theta=
\frac{\sum_{i=1}^{\ell}w_iy_i}{\sum_{i=1}^{\ell}w_i}.
\]

Это и есть оценка Надарая--Ватсона.

\begin{explanation}
Смысл вывода такой: мы ищем одно число \(\theta\), которое лучше всего описывает ответы ближайших объектов. Но ближайшие объекты учитываются с большими весами, а дальние — с малыми.
\end{explanation}

\subsection{Фиксированная и переменная ширина окна}

В случае фиксированного окна параметр \(h\) одинаков для всех объектов:
\[
h=\text{const}.
\]

При малом \(h\) модель становится более локальной и чувствительной к шуму. При большом \(h\) происходит сильное сглаживание, из-за чего модель может недообучаться.

Также рассматривается переменная ширина окна:
\[
h(z)=\rho(z,z^{(k+1)}),
\]
где \(z^{(k+1)}\) — \((k+1)\)-й ближайший сосед объекта \(z\).

В этом случае ширина окна адаптируется к плотности данных: в плотных областях окно становится меньше, а в разреженных — больше.

\begin{explanation}
Фиксированное окно использует один и тот же масштаб сглаживания везде. Переменное окно подстраивается под данные: если рядом много объектов, окно маленькое; если объектов мало, окно расширяется.
\end{explanation}

\subsection{Ядерные функции}

В работе сравниваются несколько ядерных функций.

Гауссовское ядро:
\[
K(r)=\exp\left(-\frac{r^2}{2}\right).
\]

Треугольное ядро:
\[
K(r)=(1-|r|)\mathbf{1}_{|r|\leq 1}.
\]

Ядро Епанечникова:
\[
K(r)=\frac{3}{4}(1-r^2)\mathbf{1}_{|r|\leq 1}.
\]

Биквадратное ядро:
\[
K(r)=\frac{15}{16}(1-r^2)^2\mathbf{1}_{|r|\leq 1}.
\]

\begin{explanation}
Ядро определяет, как быстро уменьшается влияние объекта при увеличении расстояния. Но на практике часто ширина окна влияет на качество сильнее, чем конкретный вид ядра.
\end{explanation}

\subsection{Подбор параметров}

Для выбора параметров модели используется кросс-валидация или leave-one-out оценка. При LOO для каждого объекта \(z_i\) модель строит прогноз без использования самого объекта \(z_i\). Затем вычисляется ошибка:
\[
LOO =
\frac{1}{\ell}
\sum_{i=1}^{\ell}
(y_i-\hat y_i^{(-i)})^2.
\]

Подбираются следующие параметры:
\begin{itemize}
    \item ширина окна \(h\) для фиксированного окна;
    \item число соседей \(k\) для переменного окна;
    \item вид ядра \(K(r)\).
\end{itemize}

Оптимальными считаются параметры, при которых ошибка на валидации минимальна.

% ============================================================
% 7. ЭКСПЕРИМЕНТЫ ПО МЕТРИЧЕСКИМ МЕТОДАМ
% ============================================================

\section{Экспериментальное исследование метрических методов}

\subsection{Одномерная иллюстрация ядерного сглаживания}

Для наглядного анализа ядерной регрессии была рассмотрена одномерная синтетическая задача:
\[
y=\sin x+\varepsilon.
\]

Такая задача позволяет визуально показать, как параметр \(h\) влияет на гладкость прогноза.

\begin{figure}[H]
    \centering
    \includegraphics[width=0.85\textwidth]{report_outputs/figures/onedim_nw_different_h.png}
    \caption{Ядерная регрессия при разных значениях ширины окна}
    \label{fig:onedim_nw_different_h}
\end{figure}

\begin{figure}[H]
    \centering
    \includegraphics[width=0.85\textwidth]{report_outputs/figures/onedim_rmse_vs_h.png}
    \caption{Зависимость ошибки от ширины окна в одномерной задаче}
    \label{fig:onedim_rmse_vs_h}
\end{figure}

\subsection{Влияние ширины окна}

Для ядерной регрессии с фиксированным окном исследовалась зависимость качества от параметра \(h\). Малые значения \(h\) приводят к высокой чувствительности к шуму, а слишком большие значения — к чрезмерному сглаживанию.

\begin{figure}[H]
    \centering
    \includegraphics[width=0.85\textwidth]{report_outputs/figures/kernel_rmse_vs_h.png}
    \caption{Зависимость ошибки ядерной регрессии от ширины окна}
    \label{fig:kernel_rmse_vs_h}
\end{figure}

\subsection{Влияние числа соседей}

Для модели с переменным окном исследовалась зависимость качества от числа соседей \(k\). Малое значение \(k\) делает модель слишком локальной, а большое значение \(k\) усиливает сглаживание.

\begin{figure}[H]
    \centering
    \includegraphics[width=0.85\textwidth]{report_outputs/figures/kernel_rmse_vs_k.png}
    \caption{Зависимость ошибки от числа соседей для переменного окна}
    \label{fig:kernel_rmse_vs_k}
\end{figure}

\subsection{Сравнение ядер}

Для оценки влияния формы ядра были сравнены различные ядерные функции. Для каждого ядра подбирались оптимальные параметры, после чего сравнивались значения ошибки на тестовой выборке.

\begin{figure}[H]
    \centering
    \includegraphics[width=0.8\textwidth]{report_outputs/figures/kernel_comparison_best_rmse.png}
    \caption{Сравнение ядерных функций по качеству прогноза}
    \label{fig:kernel_comparison_best_rmse}
\end{figure}

\subsection{Сравнение фиксированного и переменного окна}

Отдельно сравнивались лучшие модели с фиксированным и переменным окном. Переменное окно может выигрывать в случае неравномерной плотности данных, так как оно адаптируется к расположению объектов.

\begin{figure}[H]
    \centering
    \includegraphics[width=0.75\textwidth]{report_outputs/figures/fixed_vs_variable_window.png}
    \caption{Сравнение фиксированного и переменного окна}
    \label{fig:fixed_vs_variable_window}
\end{figure}

\subsection{Анализ прогнозов и ошибок ядерной модели}

Для лучшей ядерной модели были построены графики истинных и предсказанных значений, а также распределение остатков.

\begin{figure}[H]
    \centering
    \includegraphics[width=0.8\textwidth]{report_outputs/figures/kernel_true_vs_pred.png}
    \caption{Истинные и предсказанные значения для ядерной регрессии}
    \label{fig:kernel_true_vs_pred}
\end{figure}

\begin{figure}[H]
    \centering
    \includegraphics[width=0.8\textwidth]{report_outputs/figures/kernel_residuals_distribution.png}
    \caption{Распределение ошибок ядерной регрессии}
    \label{fig:kernel_residuals_distribution}
\end{figure}

% ============================================================
% 8. LOWESS
% ============================================================

\section{Робастная регрессия LOWESS}

\subsection{Идея метода}

LOWESS является робастной модификацией локального сглаживания. Основная идея состоит в том, что объектам с большими ошибками постепенно уменьшаются веса.

На каждой итерации вычисляются LOO-прогнозы:
\[
a_i=
\frac{
\sum_{j=1, j\neq i}^{\ell}
y_j\gamma_jK\left(
\frac{\rho(z_i,z_j)}{h(z_i)}
\right)
}{
\sum_{j=1, j\neq i}^{\ell}
\gamma_jK\left(
\frac{\rho(z_i,z_j)}{h(z_i)}
\right)
}.
\]

Затем вычисляются ошибки:
\[
\varepsilon_i=|a_i-y_i|.
\]

После этого объектные веса обновляются:
\[
\gamma_i=
\widetilde K
\left(
\frac{|a_i-y_i|}{6\operatorname{med}\{\varepsilon_i\}}
\right).
\]

Если объект имеет большую ошибку, то его вес \(\gamma_i\) уменьшается. Поэтому выбросы начинают меньше влиять на прогноз.

\begin{explanation}
Обычная ядерная регрессия может сильно страдать от выбросов, потому что выброс всё равно участвует в усреднении. LOWESS пытается обнаружить такие объекты по большим остаткам и уменьшить их влияние.
\end{explanation}

\subsection{Исследование устойчивости к выбросам}

Для проверки устойчивости моделей в данные добавлялись искусственные выбросы. После этого сравнивались обычная ядерная регрессия и LOWESS.

\begin{figure}[H]
    \centering
    \includegraphics[width=0.85\textwidth]{report_outputs/figures/outliers_effect_rmse.png}
    \caption{Влияние выбросов на качество регрессии}
    \label{fig:outliers_effect_rmse}
\end{figure}

\subsection{Финальные веса LOWESS}

После стабилизации алгоритма LOWESS можно проанализировать финальные веса объектов. Малые значения \(\gamma_i\) соответствуют объектам, которые алгоритм посчитал потенциальными выбросами.

\begin{figure}[H]
    \centering
    \includegraphics[width=0.85\textwidth]{report_outputs/figures/lowess_final_weights.png}
    \caption{Финальные веса объектов в LOWESS}
    \label{fig:lowess_final_weights}
\end{figure}

\subsection{Одномерная иллюстрация LOWESS}

Для наглядности также была рассмотрена одномерная задача с выбросами. На ней можно увидеть, как обычное сглаживание и LOWESS ведут себя при наличии аномальных наблюдений.

\begin{figure}[H]
    \centering
    \includegraphics[width=0.85\textwidth]{report_outputs/figures/onedim_lowess_outliers.png}
    \caption{Сравнение обычного сглаживания и LOWESS при выбросах}
    \label{fig:onedim_lowess_outliers}
\end{figure}

% ============================================================
% 9. ИТОГОВОЕ СРАВНЕНИЕ
% ============================================================

\section{Сравнение линейного и метрического подходов}

В работе были рассмотрены два класса методов.

Первый класс — линейные модели по функциональным признакам. Их преимущество состоит в интерпретируемости: модель можно записать через весовую функцию \(w(t)\), которая показывает вклад разных участков исходной функции в прогноз.

Второй класс — метрические методы регрессии. Они являются более гибкими, так как не предполагают глобальной линейной зависимости между признаками и ответом. Однако они сильнее зависят от выбора метрики, масштаба признаков, ширины окна и ядра.

Ridge-регрессия повышает устойчивость линейной модели за счёт штрафа на большие коэффициенты. LOWESS, в свою очередь, повышает устойчивость метрического подхода при наличии выбросов.

\begin{explanation}
Если нужна интерпретируемая модель, удобнее использовать линейный подход через \(w(t)\). Если зависимость сложная и нелинейная, могут быть полезны метрические методы. Если в данных есть выбросы, LOWESS может оказаться устойчивее обычного ядерного сглаживания.
\end{explanation}

\begin{table}[H]
\centering
\caption{Качественное сравнение рассмотренных моделей}
\begin{tabular}{|p{5.2cm}|p{4.3cm}|p{4.3cm}|}
\hline
\textbf{Метод} & \textbf{Преимущества} & \textbf{Ограничения} \\
\hline
Линейная регрессия по функциональным признакам 
& Хорошо интерпретируется через весовую функцию \(w(t)\) 
& Требует, чтобы зависимость была близка к линейной \\
\hline
Ridge-регрессия 
& Более устойчива при большом числе признаков и похожих функционалах 
& Требует подбора параметра регуляризации \(\lambda\) \\
\hline
Надарая--Ватсона с фиксированным окном 
& Простая непараметрическая модель, учитывающая близость объектов 
& Сильно зависит от выбора ширины окна \(h\) \\
\hline
Надарая--Ватсона с переменным окном 
& Адаптируется к плотности данных 
& Требует подбора числа соседей \(k\) \\
\hline
LOWESS 
& Устойчивее к выбросам за счёт перевзвешивания объектов 
& Сложнее вычислительно и требует итерационного процесса \\
\hline
\end{tabular}
\label{tab:qualitative_comparison}
\end{table}

\subsection{Смысловая интерпретация результатов}

С точки зрения функционального анализа линейная модель использует конечное число линейных функционалов для приближённого описания зависимости ответа от функции. Это означает, что вместо всей функции \(x(t)\) модель использует её проекции на выбранные направления \(\varphi_k(t)\).

Метрический подход использует уже не сами функции, а конечномерные представления, полученные через функционалы. Поэтому качество метрических методов зависит от того, насколько хорошо выбранные функционалы сохраняют существенную информацию о функциях.

Наиболее важным этапом оказывается построение признакового пространства. Если функционалы выбраны удачно, то и линейная модель, и метрические методы получают достаточно информации о форме исходной функции. Если же функционалы плохо отражают структуру данных, то дальнейшие методы работают с неполным представлением объекта.

Таким образом, общий результат можно сформулировать следующим образом: построение функциональных признаков является ключевым этапом, от которого зависит качество как линейных, так и метрических моделей.

% ============================================================
% ЗАКЛЮЧЕНИЕ
% ============================================================

\section*{Заключение}
\addcontentsline{toc}{section}{Заключение}

В работе была рассмотрена задача регрессии по функциональным данным. Исходные объекты задавались функциями на интервале \([0,1]\), после чего преобразовывались в конечномерные векторы признаков с помощью линейных функционалов.

Были рассмотрены разные системы функционалов: индикаторы подотрезков, тригонометрический базис и полиномиальный базис. На полученных признаках была построена линейная регрессионная модель, выведены нормальные уравнения метода наименьших квадратов, а также реализована ridge-регрессия. Дополнительно была восстановлена весовая функция \(w(t)\), позволяющая интерпретировать вклад разных участков исходной функции в прогноз.

Затем те же функциональные признаки были использованы для метрических методов регрессии. Были реализованы ядерная регрессия Надарая--Ватсона с фиксированной и переменной шириной окна, подбор параметров, сравнение различных ядер, а также робастная схема LOWESS.

По результатам экспериментов было показано, что линейная модель хорошо работает в случае, когда зависимость ответа от функции действительно близка к линейной комбинации функционалов. Её преимуществом является интерпретируемость, так как модель можно представить через весовую функцию \(w(t)\).

Метрические методы являются более гибкими и не требуют явного задания линейной формы зависимости. Однако их качество сильнее зависит от выбора метрики, ширины окна, ядра и масштаба признаков. При наличии выбросов робастная схема LOWESS позволяет уменьшить влияние аномальных объектов и повысить устойчивость прогноза.

Таким образом, кейс 4 и кейс 6 образуют единую исследовательскую схему: сначала функциональные данные переводятся в конечномерное признаковое пространство, затем в этом пространстве сравниваются параметрические и непараметрические методы регрессии.

\newpage

% ============================================================
% СПИСОК ИСТОЧНИКОВ
% ============================================================

\section*{Список использованных источников}
\addcontentsline{toc}{section}{Список использованных источников}

\begin{enumerate}
    \item Колмогоров А. Н., Фомин С. В. Элементы теории функций и функционального анализа. — М.: Наука, 1989.
    \item Канторович Л. В., Акилов Г. П. Функциональный анализ. — М.: Наука, 1984.
    \item Kreyszig E. Introductory Functional Analysis with Applications. — Wiley, 1989.
    \item Ramsay J. O., Silverman B. W. Functional Data Analysis. — Springer, 2005.
    \item Hastie T., Tibshirani R., Friedman J. The Elements of Statistical Learning. — Springer, 2009.
    \item Bishop C. M. Pattern Recognition and Machine Learning. — Springer, 2006.
    \item Nadaraya E. A. On estimating regression. Theory of Probability and Its Applications, 1964.
    \item Watson G. S. Smooth regression analysis. Sankhya: The Indian Journal of Statistics, 1964.
    \item Cleveland W. S. Robust locally weighted regression and smoothing scatterplots. Journal of the American Statistical Association, 1979.
\end{enumerate}

\newpage

% ============================================================
% ПРИЛОЖЕНИЯ
% ============================================================

\appendix

\section{Дополнительные материалы эксперимента}

В ходе работы были автоматически сформированы таблицы с результатами экспериментов. Они сохранены в формате CSV в папке:
\[
\texttt{report\_outputs/tables/}.
\]

К основным таблицам относятся:
\begin{itemize}
    \item \texttt{dataset\_description.csv} — описание параметров датасета;
    \item \texttt{linearity\_check.csv} — численная проверка линейности функционалов;
    \item \texttt{linear\_basis\_comparison.csv} — сравнение систем функционалов;
    \item \texttt{best\_linear\_models.csv} — лучшие линейные модели;
    \item \texttt{kernel\_fixed\_window\_results.csv} — результаты ядерной регрессии с фиксированным окном;
    \item \texttt{kernel\_variable\_window\_results.csv} — результаты ядерной регрессии с переменным окном;
    \item \texttt{best\_kernel\_models.csv} — лучшие метрические модели;
    \item \texttt{lowess\_outlier\_comparison.csv} — сравнение LOWESS и обычного сглаживания при выбросах;
    \item \texttt{lowess\_final\_weights.csv} — финальные веса объектов в LOWESS;
    \item \texttt{final\_model\_comparison.csv} — итоговое сравнение моделей.
\end{itemize}

\section{Исходный код}

Исходный код был реализован на языке Python. В ходе работы были реализованы следующие модули:
\begin{itemize}
    \item генерация синтетических функциональных данных;
    \item построение систем линейных функционалов;
    \item численное вычисление интегральных признаков;
    \item линейная регрессия и ridge-регрессия;
    \item ядерная регрессия Надарая--Ватсона;
    \item модель с переменной шириной окна;
    \item робастная схема LOWESS;
    \item построение графиков и таблиц с результатами.
\end{itemize}

Код позволяет воспроизвести все графики и таблицы, использованные в отчёте.

\end{document}